\chapter{Simulation Study — BASIS Planning Stage}
\label{chap:sim-plan}

This chapter documents the planning decisions for the Bayesian simulation study that underpins the comparison of Hamiltonian Monte Carlo (HMC) and Metropolis–Hastings (MH) estimators for weighted interval-censored log-logistic accelerated failure-time (AFT) models. The narrative follows the Planning phase of the BASIS framework \citep{Kelter2023BASIS} and is cross-referenced with ADEMP guidance for transparent simulation studies \citep{Morris2019Simulation}. Subsequent chapters on code, execution, analysis, and reporting build directly on the specifications recorded here.

\section{3.1 Overview and Alignment with BASIS}
\label{sec:basis-overview}
The simulation study targets HIV survey settings in which seroconversion times are interval censored and analyses rely on survey weights. During planning, BASIS prompts articulation of objectives (Aims), the synthetic cohort design (Design), the shared Bayesian model and prior structure (Model), execution logistics including random seeds and computing resources (Execution), and the evaluation criteria for comparing samplers (Performance). Table~\ref{tab:basis-mapping} provides the explicit mapping between BASIS components and ADEMP elements to signpost how each decision will be operationalised in later stages.

\begin{table}[htbp]
\centering
\caption{Mapping between BASIS planning elements and ADEMP components for the simulation study.}
\label{tab:basis-mapping}
\begin{tabular}{p{0.22\linewidth}p{0.52\linewidth}p{0.18\linewidth}}
\toprule
{BASIS element} & {Key planning decisions in this study} & {ADEMP link} \\
\midrule
Aims & Neutral comparison of HMC and MH for weighted interval-censored AFT models, framed by HIV survey analyses. & A \\
Design & Scenario grid spanning sample size, censoring intensity, and weight dispersion to emulate population-based surveys. & D \\
Model & Common log-logistic AFT specification with survey-weighted pseudo-likelihood and aligned priors. & E, M \\
Execution & Deterministic seeding, chain configuration, and PBS-array orchestration on the Stellenbosch HPC cluster. & M \\
Performance & Diagnostics and accuracy metrics (bias, RMSE, ESS/s, coverage) with pre-specified decision rules. & P \\
\bottomrule
\end{tabular}
\end{table}

\section{3.2 Simulation Aims and Questions}
\label{sec:aims}
The primary aim is to ascertain whether HMC, implemented through CmdStanR, yields superior sampling efficiency and estimator accuracy compared with MH implemented in JAGS when analysing weighted interval-censored AFT models representative of HIV seroconversion studies. Secondary aims focus on understanding how censoring severity, survey-weight dispersion, and sample size impact convergence diagnostics and computational burden across samplers, as well as documenting reproducibility artefacts that enable neutral reporting \citep{Kelter2023BASIS}. The central research questions are:
\begin{enumerate}
  \item Under which data-generating conditions do efficiency and accuracy advantages of HMC persist, diminish, or reverse relative to MH?
  \item How sensitive are posterior summaries and diagnostics to prior variants, survey-weight handling, and the width of observation intervals?
  \item What execution safeguards are required to achieve comparable convergence diagnostics across samplers while maintaining traceable computation logs?
\end{enumerate}

\section{3.3 Data-generating Mechanism}
\label{sec:data-mechanism}
Event times are generated from a log-logistic AFT formulation suited to interval-censored survival data \citep{ZhangSun2010}. For each individual $i$,
\begin{equation}
\log T_i = \log(\alpha) + X_{1i}\beta_1 + X_{2i}\beta_2 + \gamma^{-1}\varepsilon_i, \qquad \varepsilon_i \sim \text{Logistic}(0,1),
\label{eq:aft-struct}
\end{equation}
where $\alpha$ controls the baseline median time, $\beta_1$ and $\beta_2$ are regression effects for sex and age-at-entry covariates, and $\gamma$ governs tail behaviour. The covariate $X_{1i}$ is a Bernoulli indicator (e.g., sex), while $X_{2i}$ is derived from a scaled Beta distribution to mimic age-at-first-risk. Survey visit schedules yield case-II interval censoring through lower and upper bounds $(L_i,R_i)$; when a right endpoint is missing, $R_i$ is set to $\infty$. The combination of true event times, visit schedules, and administrative truncation generates observed data $(L_i, R_i, X_{1i}, X_{2i}, w_i)$ with status indicators consistent with the log-logistic likelihood.

Survey weights $w_i$ are sampled to emulate household-based HIV surveillance designs and are incorporated via a pseudo-likelihood to respect complex survey structure \citep{Si2020Bayesian}. Three weight regimes (none, low dispersion, high dispersion) cross with three censoring targets (10\%, 30\%, 50\%) at each of the three sample sizes ($n=200$, $n=2000$, $n=10000$), generating 27 distinct scenarios. Each scenario is populated with 200 replicates, producing 1800 datasets per sample size and 5400 datasets overall, as summarised in Table~\ref{tab:scenario-grid}.

\begin{table}[htbp]
\centering
\caption{Scenario grid defined during the BASIS planning stage.}
\label{tab:scenario-grid}
\begin{tabular}{l l l c c}
\toprule
Sample size & Censoring targets & Weight regimes & Scenarios & Replicates per scenario \\
\midrule
200 & 0.10, 0.30, 0.50 & high, low, none & 9 & 200 \\
2000 & 0.10, 0.30, 0.50 & high, low, none & 9 & 200 \\
10000 & 0.10, 0.30, 0.50 & high, low, none & 9 & 200 \\
\midrule
\multicolumn{4}{r}{Total datasets} & 5400 \\
\bottomrule
\end{tabular}
\end{table}

\section{3.4 Bayesian Model and Priors}
\label{sec:model-priors}
Both samplers target the same survey-weighted log-logistic AFT model; Stan code is compiled via CmdStanR and the MH implementation is executed through JAGS, ensuring that differences in posterior behaviour can be attributed to the sampling algorithms rather than to model misspecification. The likelihood contribution for each observation combines the survey weight with the interval-censored log-logistic cumulative distribution function \citep{AndersonBergman2024}. Different variations of model and prior parameters were explored and graphically assessed to confirm that simulated data adhered to the specified distributional assumptions. HMC settings follow best practices for NUTS implementations in Stan \citep{ThomasTu2021HMC}, while the MH specification mirrors the same parameterisation on the log scale to stabilise mixing.

\section{3.5 Planned Comparisons and Performance Metrics}
\label{sec:performance-plan}
\subsection*{Execution plan}
Each scenario is analysed with four parallel chains per sampler, using 1000 warm-up (HMC) or adaptation/burn-in (MH) iterations and 5000 post warm-up draws per chain, matching the configuration in \texttt{run\_fits.R}. Deterministic seeds anchor replication: a global seed of 2025 is applied to Stan, and JAGS chains receive offsets of $2025 +$ replicate index to maintain comparability across methods. Jobs are dispatched as PBS arrays on the Stellenbosch University HPC cluster, reserving one CPU core per chain and harmonising OpenMP and BLAS threading to avoid oversubscription. The software stack comprises R~4.5.1, \texttt{cmdstanr}~0.9.0 with CmdStan~2.36.0, and JAGS~4.3.2, with supporting R packages (\texttt{posterior}, \texttt{coda}, \texttt{future}, \texttt{future.apply}) loaded explicitly. All simulation scripts were version-controlled using Git to ensure reproducibility and traceability of code updates.

\subsection*{Performance criteria}
Performance measures align with ADEMP recommendations \citep{Morris2019Simulation}. Accuracy is summarised via Monte Carlo bias, root mean squared error, and Monte Carlo standard errors for key estimands. Interval coverage is recorded for central 95\% credible intervals with binomial Monte Carlo error. Efficiency is evaluated through bulk and tail effective sample sizes, effective samples per second, and split $\hat{R}$ statistics. Runtime, divergence counts, maximum treedepth hits, energy-based Bayesian fraction of missing information, and MH acceptance/auto-correlation diagnostics provide auxiliary convergence evidence. Replicates breaching $\hat{R} > 1.01$ or effective sample sizes below 100 are earmarked for reruns, and persistent failures are logged for reporting. Planned outputs include scenario-level tables and figures that contrast accuracy and efficiency across samplers while maintaining neutrality.

\begin{figure}[htbp]
\centering
\fbox{\parbox{0.8\linewidth}{\centering \TODO{Insert planned workflow schematic linking BASIS planning decisions to execution artefacts.}}}
\caption{Placeholder for the planning-stage workflow diagram.}
\label{fig:basis-workflow}
\end{figure}
